% Provenance (never printed): see sources/outline_sources.md (Year 2: tholeiitic, calc-alkaline
% and alkaline series; the TAS diagram; settings; isotopic signatures of mantle and crust).
% Classification from BGS RR 99-06 (after Streckeisen 1976, Le Bas et al. 1986, Le Maitre et
% al. 1989); series and data from Vermeesch and Pease (2021) and the GeoplotR package;
% curves and counts from figdata/bachelor-2/series.py; exercise and problem values in the
% Book 4 scratch calculators calc/ch09.py and calc/ch09pb.py.
\chapter{मैग्मीय श्रेणियाँ और विवर्तनिक परिवेश}\label{ch:b2:magmatic-series}

किसी शैलिकीविद को बिना मानचित्र के एक रासायनिक विश्लेषण दीजिए, और वह प्रायः बता सकती है कि
लावा किसी कटक पर फूटा, किसी महासागरीय द्वीप पर, या धँसती प्लेट के ऊपर। आइसलैंड के लावा
और पश्चिमी उत्तर अमेरिका के कैस्केड ज्वालामुखियों के लावा, दोनों
बेसाल्ट से रायोलाइट तक जाते हैं; सैकड़ों की संख्या में विश्लेषित करने पर वे विकसित होते हुए दो भिन्न पथों पर
चलते हैं, एक लोहे में निरंतर अधिक समृद्ध होता जाता है, दूसरा नहीं। यह अध्याय पहले
आग्नेय चट्टानों का नामकरण करता है, उनके खनिजों से जब वे दिखाई दें, और उनके
रसायन से जब न दिखें; फिर वह उन्हें श्रेणियों में बाँटता है, यानी विभेदन से संबंधित
चट्टानों के परिवारों में, और हर श्रेणी को एक विवर्तनिक परिवेश से जोड़ता है, उन
आरेखों के साथ जिनसे शैलिकीविद किसी विश्लेषण से परिवेश पढ़ते हैं, और उनकी सीमाओं के साथ।
% ledger: vp21:series

\begin{recall}
\textbf{पुस्तक 2 (उच्च माध्यमिक) से}: बेसाल्ट, एंडेसाइट, रायोलाइट, गैब्रो और ग्रेनाइट;
फ़ेल्सिक और मैफ़िक चट्टानें; ज्वालामुखीय और पातालीय बनावटें।
\textbf{पुस्तक 3 (स्नातक वर्ष 1) से}: मध्य-महासागरीय कटक, क्षेपण क्षेत्र और ज्वालामुखी चाप, हॉटस्पॉट और बृहत
आग्नेय प्रांत, महाद्वीपीय दरारें। \textbf{\cref{ch:b2:elements-in-the-earth} से}:
MORB और OIB, आद्य-प्रावार-प्रसामान्यीकृत (स्पाइडर) आरेख, संगत और
\omterm{def:b2:elements-in-the-earth:partition}{असंगत तत्त्व}। \textbf{\cref{ch:b2:isotope-geochemistry} से}: क्षीण
प्रावार और समृद्ध समस्थानिक भंडार। \textbf{\cref{ch:b2:magmatic-differentiation} से}:
\omterm{def:b2:magmatic-differentiation:crystallisation}{भिन्नात्मक क्रिस्टलन}, \omterm{def:b2:magmatic-differentiation:lld}{द्रव का
अवरोहण-पथ}, \omterm{def:b2:magmatic-differentiation:harker}{हार्कर आरेख} और रैले नियम।
\end{recall}

\section{आग्नेय चट्टानों का नामकरण}\label{sec:b2:magmatic-series:naming}

स्थूल-कणिक चट्टान अपने खनिज दिखाती है, और उन्हीं से उसका नाम रखा जाता है। लावा प्रायः इसके लिए बहुत
सूक्ष्म-कणिक या बहुत काचीय होता है, और उसका नाम उसके रासायनिक विश्लेषण से रखा जाता है। दोनों
पद्धतियाँ अंतरराष्ट्रीय परिपाटियाँ हैं: वे वर्णन करती हैं, और इस बारे में कुछ नहीं कहतीं कि
चट्टान कैसे बनी।
% ledger: bgs:tas-use

\begin{definition}[मोडल संघटन; QAPF आरेख]\label{def:b2:magmatic-series:qapf}
किसी चट्टान का \emph{मोडल संघटन}\index{मोडल संघटन} उसके हर खनिज का आयतन के अनुसार
अनुपात है। \emph{QAPF आरेख}\index{QAPF आरेख}
\qty{90}{\percent} से कम मैफ़िक खनिजों वाली स्थूल-कणिक आग्नेय चट्टान का नाम
चार मोडल घटकों से रखता है, जिन्हें 100 के योग पर पुनः परिकलित किया जाता है: क्वार्ट्ज़ (Q), क्षार-फ़ेल्डस्पार
(A), प्लैजियोक्लेज़ (P) और फ़ेल्डस्पैथॉइड (F)। चूँकि क्वार्ट्ज़ और फ़ेल्डस्पैथॉइड एक साथ
क्रिस्टलित नहीं होते, हर चट्टान दो त्रिभुजों, QAP या FAP, में से एक में पड़ती है, जो
अपनी उभयनिष्ठ भुजा AP के अनुदिश जुड़े हैं।
\end{definition}
% ledger: bgs:qapf

\begin{omfigure}
\centering
\begin{tikzpicture}[font=\scriptsize, x=1cm, y=1cm]
\def\L{6.4}
\pgfmathsetmacro{\H}{\L*0.8660254}
% upper point (q, p): q = Q fraction, p = P/(A+P)
\newcommand{\qapfup}[2]{({(#1)*\L/2 + (#2)*(1-(#1))*\L}, {(#1)*\H})}
\newcommand{\dn}[2]{({(#1)*\L/2 + (#2)*(1-(#1))*\L}, {-(#1)*\H})}
\draw[thick] (0,0) -- (\L/2,\H) -- (\L,0) -- (\L/2,-\H) -- cycle;
\draw (0,0) -- (\L,0);
\foreach \q in {0.9,0.6,0.2,0.05} \draw \qapfup{\q}{0} -- \qapfup{\q}{1};
\foreach \q in {0.1,0.6} \draw \dn{\q}{0} -- \dn{\q}{1};
\foreach \p in {0.1,0.65,0.9} \draw \qapfup{0.6}{\p} -- \qapfup{0}{\p};
\draw[dashed] \qapfup{0.6}{0.35} -- \qapfup{0.2}{0.35};
\draw \qapfup{0.2}{0.35} -- \qapfup{0}{0.35};
\foreach \p in {0.1,0.35,0.65,0.9} \draw \dn{0}{\p} -- \dn{0.1}{\p};
\foreach \p in {0.1,0.5,0.9} \draw \dn{0.1}{\p} -- \dn{0.6}{\p};
\node[anchor=south] at (\L/2,\H) {\normalsize Q};
\node[anchor=east] at (0,0) {\normalsize A};
\node[anchor=west] at (\L,0) {\normalsize P};
\node[anchor=north] at (\L/2,-\H) {\normalsize F};
\node at \qapfup{0.94}{0.5} {1a};
\node at \qapfup{0.75}{0.5} {1b};
\node[font=\tiny] at \qapfup{0.4}{0.05} {2};
\node at \qapfup{0.4}{0.225} {3a};
\node at \qapfup{0.4}{0.5} {3b};
\node at \qapfup{0.4}{0.775} {4};
\node[font=\tiny] at \qapfup{0.4}{0.95} {5};
\foreach \p/\n in {0.225/7*,0.5/8*,0.775/9*} \node at \qapfup{0.125}{\p} {\n};
\node[font=\tiny] at \qapfup{0.125}{0.05} {6*};
\node[font=\tiny] at \qapfup{0.125}{0.95} {10*};
\foreach \p/\n in {0.225/7,0.5/8,0.775/9} \node at \qapfup{0.025}{\p} {\n};
\foreach \p/\n in {0.225/7',0.5/8',0.775/9'} \node at \dn{0.05}{\p} {\n};
\node[font=\tiny] at \dn{0.3}{0.05} {11};
\node at \dn{0.35}{0.3} {12};
\node at \dn{0.35}{0.7} {13};
\node[font=\tiny] at \dn{0.3}{0.95} {14};
\node at \dn{0.75}{0.5} {15};
\node[anchor=east] at \qapfup{0.9}{0} {90};
\node[anchor=east] at \qapfup{0.6}{0} {60};
\node[anchor=east] at \qapfup{0.2}{0} {20};
\node[anchor=east] at \dn{0.1}{0} {10};
\node[anchor=east] at \dn{0.6}{0} {60};
\foreach \p/\n in {0.1/10,0.35/35,0.65/65,0.9/90} \node[anchor=south, fill=white, inner sep=0.5pt] at \qapfup{0.6}{\p} {\tiny\n};
\node[anchor=west, align=left, font=\scriptsize] at (\L+0.9,1.6) {ऊपरी त्रिभुज: Q + A + P\\में Q का \%;\\
  अचर $100\,P/(A+P)$\\की रेखाएँ\\10, 35, 65, 90 पर};
\node[anchor=west, align=left, font=\scriptsize] at (\L+0.9,-1.7) {निचला त्रिभुज: F + A + P\\में F का \%;\\
  $100\,P/(A+P)$\\10, 50, 90 पर};
\end{tikzpicture}
\par\smallskip{\small स्थूल-कणिक आग्नेय चट्टानों का QAPF वर्गीकरण। 1a
क्वार्ट्ज़ोलाइट, 1b क्वार्ट्ज़-समृद्ध ग्रेनिटॉइड, 2 क्षार-फ़ेल्डस्पार ग्रेनाइट, 3a साइनोग्रेनाइट और
3b मॉन्ज़ोग्रेनाइट (मिलकर ग्रेनाइट), 4 ग्रैनोडायोराइट, 5 टोनालाइट; 6*--10*
6--10 की क्वार्ट्ज़-युक्त किस्में: 6 क्षार-फ़ेल्डस्पार साइनाइट, 7 साइनाइट, 8 मॉन्ज़ोनाइट, 9 मॉन्ज़ोडायोराइट
या मॉन्ज़ोगैब्रो, 10 डायोराइट, गैब्रो या ऐनॉर्थोसाइट; 6'--10' उनकी फ़ॉइड-युक्त
किस्में; 11 फ़ॉइड साइनाइट, 12 फ़ॉइड मॉन्ज़ोसाइनाइट, 13 फ़ॉइड मॉन्ज़ोडायोराइट या
मॉन्ज़ोगैब्रो, 14 फ़ॉइड डायोराइट या गैब्रो, 15 फ़ॉइडोलाइट। किनारों के अनुदिश पतले क्षेत्र 6, 6', 10 और 10'
क्रमांकित नहीं हैं। ब्रिटिश भूवैज्ञानिक सर्वेक्षण की पद्धति के अनुसार (जो
स्ट्रेकाइज़न 1976 पर आधारित है)।}
\end{omfigure}
% ledger: bgs:qapf

क्षेत्र 10 में तीन नाम हैं। \omterm{def:b2:magmatic-series:qapf}{QAPF आरेख} डायोराइट, गैब्रो और
ऐनॉर्थोसाइट को अलग नहीं करता, न ही गैब्रोई चट्टानों को आपस में; यह काम आगे के आरेख करते हैं,
जो प्लैजियोक्लेज़ और मैफ़िक खनिजों पर आधारित हैं।

\begin{definition}[कुल क्षार--सिलिका आरेख]\label{def:b2:magmatic-series:tas}
\emph{कुल क्षार--सिलिका आरेख}\index{कुल क्षार--सिलिका आरेख} (TAS
आरेख) ऐसी ज्वालामुखीय चट्टान का नाम, जिसका मोड मापा नहीं जा सकता,
\ce{Na2O + K2O} के द्रव्यमान-प्रतिशत से \ce{SiO2} के विरुद्ध रखता है, जब विश्लेषण को
जल और कार्बन डाइऑक्साइड से मुक्त \qty{100}{\percent} पर पुनः परिकलित कर लिया गया हो। यह
सीधी रेखाओं द्वारा पंद्रह क्षेत्रों में बँटा है, जिनके प्रतिच्छेदों के निर्देशांक निश्चित हैं
(अगले अनुच्छेद के बाद का चित्र)।
\end{definition}
% ledger: bgs:tas, bgs:tas-use

आरेख में क्षेत्रों के तीन समूह हैं। नीचे के अनुदिश, कम क्षारों पर,
सामान्य श्रेणी के नाम हैं, पिक्रोबेसाल्ट और बेसाल्ट से रायोलाइट तक, जो 45, 52,
57 और \qty{63}{\percent} सिलिका पर, और रायोलाइट के लिए एक ढलवाँ रेखा से, अलग होते हैं। उनके ऊपर S क्षेत्र और ट्रैकाइट
उसी सिलिका पर क्षारों में अधिक समृद्ध चट्टानें रखते हैं; U क्षेत्र, फ़ोनोलाइट और फ़ॉइडाइट ऐसी चट्टानें रखते हैं जो
क्षारों में इतनी समृद्ध और सिलिका में इतनी निर्धन हैं कि वे फ़ेल्डस्पैथॉइड क्रिस्टलित करती हैं। कुछ क्षेत्र
एक दूसरी कसौटी से उपविभाजित हैं: ट्रैकीबेसाल्ट को हवाइआइट कहते हैं यदि $\ce{Na2O} - 2
\ge \ce{K2O}$, अन्यथा पोटैशियमी ट्रैकीबेसाल्ट; जिस टेफ़्राइट के नॉर्म (नीचे) में
\qty{10}{\percent} से अधिक ओलिवीन हो वह बेसानाइट है; जिस रायोलाइट या ट्रैकाइट का
आणविक $(\ce{Na2O + K2O})/\ce{Al2O3}$ 1 से अधिक हो उसे अतिक्षारीय कहते हैं। यह पद्धति
\qty{2}{\percent} से कम जल और \qty{0.5}{\percent} से कम
कार्बन डाइऑक्साइड वाले विश्लेषण माँगती है: परिवर्तन जल जोड़ता है और क्षारों को खिसकाता है, और \omterm{def:b2:magmatic-differentiation:cumulate}{क्यूम्यूलेट}
द्रव नहीं है।
% ledger: bgs:tas, bgs:tas-use

\begin{omfigure}
\centering
\begin{tikzpicture}
\begin{axis}[omrecord, width=0.86\linewidth, height=8.6cm,
  xmin=37, xmax=80, ymin=0, ymax=16, xtick={37,41,...,77}, ytick={0,2,...,16},
  xlabel={\ce{SiO2} (\unit{\percent})}, ylabel={\ce{Na2O + K2O} (\unit{\percent})},
  legend style={font=\scriptsize, at={(0.98,0.98)}, anchor=north east},
  tick label style={font=\scriptsize}, label style={font=\small}]
\addplot[only marks, mark=*, mark size=0.55pt, blue!70!black, opacity=0.6] table[x=SiO2, y=alk] {figdata/out/bachelor-2-series-cath-th.dat};
\addlegendentry{आइसलैंड}
\addplot[only marks, mark=*, mark size=0.55pt, red!75!black, opacity=0.6] table[x=SiO2, y=alk] {figdata/out/bachelor-2-series-cath-ca.dat};
\addlegendentry{कैस्केड}
\foreach \f in {Pc,B,O1,O2,O3,R,S1,S2,S3,T,U1,U2,U3,Ph,F}
  \addplot[black, thin, forget plot] table[x=SiO2, y=alk] {figdata/out/bachelor-2-series-tas-\f.dat};
\node[font=\tiny] at (axis cs:43,1.6) {Pc};
\node[font=\tiny, fill=white, inner sep=0.5pt] at (axis cs:50.5,0.8) {B};
\node[font=\tiny] at (axis cs:54.5,2.2) {O1};
\node[font=\tiny] at (axis cs:60,2.2) {O2};
\node[font=\tiny] at (axis cs:66.5,2.2) {O3};
\node[font=\tiny] at (axis cs:74.5,11) {R};
\node[font=\tiny, fill=white, inner sep=0.5pt] at (axis cs:49,5.8) {S1};
\node[font=\tiny] at (axis cs:52.5,7.4) {S2};
\node[font=\tiny] at (axis cs:57.5,8.9) {S3};
\node[font=\tiny] at (axis cs:64.5,12) {T};
\node[font=\tiny] at (axis cs:44,6.6) {U1};
\node[font=\tiny] at (axis cs:48.7,9.4) {U2};
\node[font=\tiny] at (axis cs:52.5,11.8) {U3};
\node[font=\tiny] at (axis cs:57,14.6) {Ph};
\node[font=\tiny] at (axis cs:39,10) {F};
\end{axis}
\end{tikzpicture}
\par\smallskip{\small TAS आरेख (ले बा और अन्य, 1986, के अनुसार, जैसा
ब्रिटिश भूवैज्ञानिक सर्वेक्षण देता है), आइसलैंड के 456 लावाओं और अंतर्वेधी चट्टानों (नीला) और
कैस्केड के 580 (लाल) के साथ, ऑक्साइड विश्लेषित रूप में। Pc पिक्रोबेसाल्ट, B बेसाल्ट, O1 बेसाल्टी
एंडेसाइट, O2 एंडेसाइट, O3 डेसाइट, R रायोलाइट; S1 ट्रैकीबेसाल्ट, S2 बेसाल्टी
ट्रैकीएंडेसाइट, S3 ट्रैकीएंडेसाइट, T ट्रैकाइट या ट्रैकीडेसाइट; U1 टेफ़्राइट या
बेसानाइट, U2 फ़ोनोटेफ़्राइट, U3 टेफ़्रीफ़ोनोलाइट, Ph फ़ोनोलाइट, F फ़ॉइडाइट।}
\end{omfigure}
% ledger: bgs:tas, geoplotr:tas, data:b4-cath


\begin{definition}[सिलिका संतृप्ति]\label{def:b2:magmatic-series:saturation}
कोई \omterm{def:b2:magma-emplacement:magma}{मैग्मा} या चट्टान \emph{सिलिका-अतिसंतृप्त}\index{सिलिका संतृप्ति} होती है यदि उसमें
उससे अधिक सिलिका हो जितनी उसके अन्य ऑक्साइड सिलिकेट खनिजों में जोड़ सकते हैं, जिससे
क्वार्ट्ज़ (या कोई अन्य सिलिका खनिज) प्रकट होता है; \emph{सिलिका-असंतृप्त} यदि उसमें इतनी
कम सिलिका हो कि उसके सारे क्षार फ़ेल्डस्पार में न बदल सकें, जिससे उसके स्थान पर कोई फ़ेल्डस्पैथॉइड (नेफ़लीन,
ल्यूसाइट) प्रकट होता है, जो क्वार्ट्ज़ के साथ कभी नहीं आता; और \emph{सिलिका-संतृप्त}
इन दोनों के बीच।
\end{definition}

जिस लावा के खनिज गिने नहीं जा सकते, उसकी संतृप्ति एक \emph{नॉर्म} पर आँकी जाती है: एक
मानक परिकलन जो विश्लेषण के ऑक्साइडों को एक निश्चित क्रम में आदर्श खनिजों में बाँटता है,
पहले क्षारों को फ़ेल्डस्पार में, और अतिरिक्त सिलिका को
क्वार्ट्ज़ के रूप में या कमी को फ़ेल्डस्पैथॉइड के रूप में दर्ज करता है। प्रचलित नॉर्म, CIPW नॉर्म,
बेसाल्ट के दो प्रकारों को अलग करता है: जिस बेसाल्ट के नॉर्म में नेफ़लीन हो वह क्षारीय बेसाल्ट है; जिसमें
न हो वह उपक्षारीय बेसाल्ट है, और इनमें से थोलिआइटी बेसाल्टों के नॉर्म में
हाइपरस्थीन (ऑर्थोपाइरॉक्सीन) होता है और, प्रायः, मोड में निम्न-कैल्शियम पाइरॉक्सीन।
\Cref{exo:b2:magmatic-series:11} इसका एक सरलीकृत रूप चलाता है।
% ledger: bgs:tas-use, bgs:series-terms

\section{मैग्मीय श्रेणियाँ}\label{sec:b2:magmatic-series:series}

एक ज्वालामुखी, या एक प्रांत, के लावा शायद ही कभी एक जैसे होते हैं, पर वे संबंधित होते हैं: वे
समान जनक मैग्माओं से \cref{ch:b2:magmatic-differentiation} के विभेदन द्वारा
व्युत्पन्न होते हैं, और साथ मिलकर एक अवरोहण-पथ खींचते हैं।

\begin{definition}[मैग्मीय श्रेणी]\label{def:b2:magmatic-series:series}
\emph{मैग्मीय श्रेणी}\index{मैग्मीय श्रेणी} किसी एक
मैग्मीय चक्र या प्रांत की आग्नेय चट्टानों का ऐसा समूह है जो समान जनक मैग्माओं से विभेदन द्वारा संबंधित हो और
विशिष्ट रासायनिक तथा खनिजीय लक्षण साझा करता हो। सामान्यतः तीन
पहचानी जाती हैं। \emph{थोलिआइटी श्रेणी}\index{थोलिआइटी श्रेणी} सिलिका में संतृप्त या थोड़े अतिसंतृप्त,
नॉर्म में हाइपरस्थीन वाले बेसाल्टों से शुरू होती है,
और अपने विभेदन के आरंभ में लोहे में स्पष्ट संवर्धन दिखाती है।
\emph{कैल्क-क्षारीय श्रेणी}\index{कैल्क-क्षारीय श्रेणी} बेसाल्ट से
एंडेसाइट होते हुए डेसाइट और रायोलाइट तक जाती है, लोहे में बहुत कम या बिना संवर्धन के, और
अपने विकसित सदस्यों में जलयुक्त खनिजों (हॉर्नब्लेंड, बायोटाइट) के साथ। \emph{क्षारीय
श्रेणी}\index{क्षारीय श्रेणी} नेफ़लीन-नॉर्मेटिव बेसाल्टों से शुरू होती है और
TAS आरेख के S और U क्षेत्रों से होते हुए ट्रैकाइट और फ़ोनोलाइट की ओर विकसित होती है।
\end{definition}
% ledger: bgs:series-terms, bgs:tas-use

ये नाम समूहों का वर्णन करते हैं, अकेली चट्टानों का नहीं: एक अकेला बेसाल्ट इनमें से किसी का भी
हो सकता है, और यही नाम भिन्न अर्थों में प्रयुक्त हुए हैं, इसलिए शैलिकीविद को
बताना चाहिए कि उनसे उसका क्या आशय है। थोलिआइटी और
\omterm{def:b2:magmatic-series:series}{कैल्क-क्षारीय श्रेणियों} का भेद लोहे पर टिका है, और जो आरेख इसे सबसे अच्छी तरह दिखाता है वह ऑक्साइडों के तीन
समूहों को आलेखित करता है।
% ledger: bgs:series-terms

\begin{definition}[AFM आरेख]\label{def:b2:magmatic-series:afm}
\emph{AFM आरेख}\index{AFM आरेख} एक त्रिभुजाकार आरेख है जिसके शीर्ष
$A = \ce{Na2O + K2O}$, $F = \ce{FeO_T}$ (सारा लोहा \ce{FeO} के रूप में) और $M = \ce{MgO}$ हैं,
द्रव्यमान-प्रतिशत में, इस प्रकार पुनः परिकलित कि $A + F + M = 100$। उस पर खींची गई एक रेखा
F की ओर के थोलिआइटी क्षेत्र को कैल्क-क्षारीय क्षेत्र से अलग करती है।
\end{definition}
% ledger: vp21:series, geoplotr:afm

विभेदित होता बेसाल्ट मैग्नीशियम को किसी भी अन्य चीज़ से तेज़ी से खोता है, क्योंकि ओलिवीन और
पाइरॉक्सीन द्रव से अधिक मैग्नीशियम-समृद्ध हैं, और अपने क्षारों को रखे रहता है, जो
असंगत हैं: उसका अवरोहण-पथ M से दूर, A की ओर, जाता है। लोहे का क्या होता है,
यही श्रेणी तय करता है। थोलिआइटी \omterm{def:b2:magma-emplacement:magma}{मैग्मा} में लोहा केवल
मैग्नीशियम-समृद्ध फ़ेरस सिलिकेटों के साथ, मैग्नीशियम से धीमे, निकलता है, जिससे अनुपात
$F/M$ बढ़ता है और प्रवृत्ति पहले F की ओर चढ़ती है, फिर A की ओर मुड़ती है: यह
\emph{फ़ेनर प्रवृत्ति} है। अधिक ऑक्सीकृत कैल्क-क्षारीय \omterm{def:b2:magma-emplacement:magma}{मैग्मा} में लोहा
मैग्नेटाइट में भी निकलता है, जिसमें फ़ेरिक लोहा होता है, और $F/M$ लगभग स्थिर रहता है, जबकि द्रव
सीधे A की ओर बढ़ता है: यह \emph{बोवेन प्रवृत्ति} है। नीचे का चित्र दोनों को
वास्तविक चट्टानों के साथ दिखाता है, आइसलैंड के थोलिआइट और कैस्केड के कैल्क-क्षारीय लावा।
1971 में आँख से खींची गई विभाजक रेखा उन्हें वहाँ अच्छी तरह अलग करती है जहाँ वे बेसाल्टी हैं;
अधिक क्षारों पर दोनों प्रवृत्तियाँ अभिसरित होती हैं और रेखा बहुत कम बताती है।
% ledger: vp21:series, geoplotr:afm, data:b4-cath

\begin{omfigure}
\centering
\begin{tikzpicture}
\begin{axis}[name=afm, width=0.5\linewidth, height=0.5\linewidth, axis equal image,
  hide axis, xmin=-0.04, xmax=1.04, ymin=-0.06, ymax=0.92, clip=false]
\addplot[only marks, mark=*, mark size=0.5pt, blue!70!black, opacity=0.6] table[x=afm_x, y=afm_y] {figdata/out/bachelor-2-series-cath-th.dat};
\addplot[only marks, mark=*, mark size=0.5pt, red!75!black, opacity=0.6] table[x=afm_x, y=afm_y] {figdata/out/bachelor-2-series-cath-ca.dat};
\addplot[very thick, black] table[x=x, y=y] {figdata/out/bachelor-2-series-afm-line.dat};
\draw[thick] (axis cs:0,0) -- (axis cs:1,0) -- (axis cs:0.5,0.8660254) -- cycle;
\node[anchor=south] at (axis cs:0.5,0.866) {F};
\node[anchor=north east] at (axis cs:0,0) {A};
\node[anchor=north west] at (axis cs:1,0) {M};
\node[font=\scriptsize, fill=white, inner sep=1pt] at (axis cs:0.52,0.62) {थोलिआइटी};
\node[font=\scriptsize, fill=white, inner sep=1pt] at (axis cs:0.33,0.12) {कैल्क-क्षारीय};
\end{axis}
\begin{axis}[omrecord, at={(afm.east)}, anchor=west, xshift=1.3cm,
  width=0.46\linewidth, height=0.46\linewidth,
  xmin=-2.5, xmax=2.2, ymin=-4.5, ymax=1.5,
  xlabel={$\ln(A/F)$}, ylabel={$\ln(M/F)$},
  legend style={font=\scriptsize, at={(0.98,0.98)}, anchor=north east},
  tick label style={font=\scriptsize}, label style={font=\small}]
\addplot[only marks, mark=*, mark size=0.5pt, blue!70!black, opacity=0.6] table[x=lnAF, y=lnMF] {figdata/out/bachelor-2-series-lr-th.dat};
\addlegendentry{आइसलैंड}
\addplot[only marks, mark=*, mark size=0.5pt, red!75!black, opacity=0.6] table[x=lnAF, y=lnMF] {figdata/out/bachelor-2-series-lr-ca.dat};
\addlegendentry{कैस्केड}
\addplot[thick, blue!50!black, forget plot] table[x=x, y=y_th] {figdata/out/bachelor-2-series-lr-fits.dat};
\addplot[thick, red!50!black, forget plot] table[x=x, y=y_ca] {figdata/out/bachelor-2-series-lr-fits.dat};
\end{axis}
\end{tikzpicture}
\par\smallskip{\small बाएँ: TAS चित्र की आइसलैंडी (नीला) और कैस्केड (लाल)
चट्टानों का \omterm{def:b2:magmatic-series:afm}{AFM आरेख}, इर्विन और बाराघर (1971) की विभाजक रेखा के साथ। दाएँ:
वही चट्टानें लघुगणकीय-अनुपात निर्देशांकों में, जहाँ हर श्रेणी लगभग सीधी
पट्टी बन जाती है; रेखाएँ लांबिक समाश्रयण हैं, जिनकी ढलानें $-1.10$ (आइसलैंड) और $-0.72$
(कैस्केड) हैं, और जो $(-3.84, 2.12)$ पर मिलती हैं, जैसा फ़र्मेश और पीज़ (2021) ने पाया।}
\end{omfigure}
% ledger: geoplotr:afm, data:b4-cath, vp21:model

विभाजक रेखा के परास के भीतर की आइसलैंडी चट्टानों में से 351
थोलिआइटी ओर पड़ती हैं और 23 दूसरी ओर; कैस्केड की चट्टानों में से 417 कैल्क-क्षारीय
ओर और 57 दूसरी ओर। वक्र एक भद्दी सीमा है, और एक सरल मॉडल
प्रवृत्तियों के आकार की व्याख्या करता है। मान लीजिए कि कोई \omterm{def:b2:magma-emplacement:magma}{मैग्मा} अपने क्षार, लोहा और मैग्नीशियम
अपनी धारित मात्राओं के समानुपाती दरों से खोता है, $\mathrm dX/\mathrm dt = -\lambda_X X$।

\begin{proposition}[लघुगणकीय-अनुपात प्रवृत्तियाँ]\label{prop:b2:magmatic-series:logratio}
यदि $A$, $F$ और $M$ स्थिरांकों $\lambda_A$, $\lambda_F$ और
$\lambda_M$ के साथ चरघातांकी रूप से क्षयित होते हैं, तो \omterm{def:b2:magma-emplacement:magma}{मैग्मा} $\ln(A/F)$ और
$\ln(M/F)$ के तल में एक सीधी रेखा पर चलता है, जिसकी ढलान $(\lambda_F - \lambda_M)/(\lambda_F - \lambda_A)$ है, जबकि
AFM त्रिभुज पर वह एक वक्र के अनुदिश चलता है।
\end{proposition}

\begin{proof}
$X = X_0 e^{-\lambda_X t}$ के साथ, $\ln(A/F) = \ln(A_0/F_0) + (\lambda_F - \lambda_A)t$
और $\ln(M/F) = \ln(M_0/F_0) + (\lambda_F - \lambda_M)t$: दोनों $t$ में रैखिक हैं, अतः
एक दूसरे में रैखिक है, बताई गई ढलान के साथ। अनुपात $A/F$ और $M/F$ त्रिभुज पर
बिंदु निश्चित करते हैं, जिसके निर्देशांक $A/(A + F + M)$, \dots\ $t$ में रैखिक नहीं हैं।
\end{proof}

आइसलैंडी चट्टानों की अधिक तीव्र ढलान का अर्थ है छोटा $\lambda_F$: लोहा
थोलिआइटी \omterm{def:b2:magma-emplacement:magma}{मैग्मा} से अधिक धीरे निकलता है (\cref{exo:b2:magmatic-series:10})। फ़र्मेश और पीज़ ने दिखाया कि ऊपर के चित्र की दोनों
पट्टियाँ एक बिंदु पर मिलती हैं, जिसे उन्होंने एक काल्पनिक साझा
स्रोत के रूप में पढ़ा।
% ledger: vp21:model

थोलिआइटी और \omterm{def:b2:magmatic-series:series}{कैल्क-क्षारीय श्रेणियाँ} अपने पोटाश में भी भिन्न हैं।
\ce{K2O}--\ce{SiO2} आरेख \qty{48}{\percent} से अधिक सिलिका वाले बेसाल्टों, एंडेसाइटों, डेसाइटों और रायोलाइटों को
दो सीधी रेखाओं द्वारा निम्न-K, मध्यम-K और उच्च-K प्रकारों में बाँटता है
(नीचे का चित्र)। यहाँ दोनों समूहों का अधिकांश भाग मध्यम-K है: यह
आरेख बताता है कि कोई श्रेणी कितना पोटाश रखती है, यह नहीं कि वह कौन-सी श्रेणी है।
% ledger: bgs:k2o-sio2, data:b4-cath

\begin{omfigure}
\centering
\begin{tikzpicture}
\begin{axis}[omrecord, width=0.78\linewidth, height=6.2cm,
  xmin=45, xmax=79, ymin=0, ymax=5,
  xlabel={\ce{SiO2} (\unit{\percent})}, ylabel={\ce{K2O} (\unit{\percent})},
  legend style={font=\scriptsize, at={(0.02,0.98)}, anchor=north west},
  tick label style={font=\scriptsize}, label style={font=\small}]
\addplot[only marks, mark=*, mark size=0.55pt, blue!70!black, opacity=0.6] table[x=SiO2, y=K2O] {figdata/out/bachelor-2-series-cath-th.dat};
\addlegendentry{आइसलैंड}
\addplot[only marks, mark=*, mark size=0.55pt, red!75!black, opacity=0.6] table[x=SiO2, y=K2O] {figdata/out/bachelor-2-series-cath-ca.dat};
\addlegendentry{कैस्केड}
\addplot[thick, black, forget plot] table[x=x, y=y_low] {figdata/out/bachelor-2-series-k2o-lines.dat};
\addplot[thick, black, forget plot] table[x=x, y=y_high] {figdata/out/bachelor-2-series-k2o-lines.dat};
\draw[thick] (axis cs:48,0) -- (axis cs:48,0.3);
\node[font=\scriptsize, fill=white, inner sep=1pt] at (axis cs:60,3.7) {उच्च-K};
\node[font=\scriptsize, fill=white, inner sep=1pt] at (axis cs:76.5,2.2) {मध्यम-K};
\node[font=\scriptsize, fill=white, inner sep=1pt] at (axis cs:60,0.35) {निम्न-K};
\end{axis}
\end{tikzpicture}
\par\smallskip{\small आइसलैंडी और कैस्केड चट्टानों के लिए सिलिका के विरुद्ध पोटाश, उन
रेखाओं के साथ जो (48, 0.3)--(68, 1.2) और (48, 1.2)--(68, 3.1) से होकर जाती हैं और निम्न-K,
मध्यम-K और उच्च-K प्रकारों को बाँटती हैं (पेचेरिलो और टेलर, 1976, के अनुसार, जैसा
अंतरराष्ट्रीय पद्धति ने अपनाया है), यहाँ 68 से आगे बढ़ाई गईं।}
\end{omfigure}
% ledger: bgs:k2o-sio2, data:b4-cath

ऑक्सीकरण भिन्न क्यों होता है? कैल्क-क्षारीय \omterm{def:b2:magma-emplacement:magma}{मैग्मा} क्षेपण क्षेत्रों के ऊपर बनते हैं, जहाँ
निर्जलित होता स्लैब प्रावार-कील को ऑक्सीकृत स्पीशीज़ देता है; ये \omterm{def:b2:magma-emplacement:magma}{मैग्मा}
मैग्नेटाइट क्रिस्टलित करते हैं, और उनके विकसित सदस्य हॉर्नब्लेंड या बायोटाइट। थोलिआइटी
\omterm{def:b2:magma-emplacement:magma}{मैग्मा} कटकों पर बनते हैं, जहाँ प्रावार विसंपीडन से पिघलता है
(\cref{ch:b2:mantle-melting}), और अधिक अपचयित होते हैं।
% ledger: vp21:series

\section{श्रेणियाँ और परिवेश}\label{sec:b2:magmatic-series:settings}

प्लेट विवर्तनिकी का हर परिवेश अपने \omterm{def:b2:magma-emplacement:magma}{मैग्मा} बनाता है, क्योंकि हर एक
प्रावार को अपने ढंग से पिघलाता है: कटकों और हॉटस्पॉटों के नीचे विसंपीडन से, क्षेपण क्षेत्रों के ऊपर
जल जोड़कर (\cref{ch:b2:mantle-melting})।

\emph{मध्य-महासागरीय कटक} थोलिआइटी बेसाल्ट उगलते हैं, यानी
\cref{ch:b2:elements-in-the-earth} के MORB, जो \omterm{def:b2:elements-in-the-earth:partition}{असंगत तत्त्वों} में निर्धन हैं और जिनमें
अपने स्रोत का क्षीण समस्थानिक चिह्न होता है। आइसलैंड, जो किसी हॉटस्पॉट के ऊपर स्थित कटक है, इस सूची में
अपने थोलिआइट जोड़ता है, रायोलाइट तक के पूरे समूह के साथ।
% ledger: vp21:series

\emph{महासागरीय द्वीप} चरणों में बढ़ते हैं। हवाई का कोई ज्वालामुखी युवा अवस्था में क्षारीय लावा उगलता है,
फिर थोलिआइटी बेसाल्ट की अपनी ढाल-संरचना बनाता है, जो उसके अंतिम आयतन का अनुमानित 80 से
\qty{95}{\percent} है, अपने पहले 1 से 2 मिलियन वर्षों के भीतर; जैसे-जैसे
वह हॉटस्पॉट से दूर खिसकता है उसके लावा फिर क्षारीय हो जाते हैं, और अंत में
विरल क्षारीय उद्गारों की एक लंबी पूँछ आती है। कम गलन अधिक क्षारीय \omterm{def:b2:magma-emplacement:magma}{मैग्मा} देता है: क्षार,
असंगत होने के कारण, छोटे \omterm{def:b2:mantle-melting:fraction}{गलन-अंश} में कम तनु होते हैं। समस्थानिक रूप से, महासागरीय-द्वीप
बेसाल्ट \cref{ch:b2:isotope-geochemistry} के समृद्ध भंडारों का नमूना लेते हैं।
% ledger: usgs14:hawaii-stages

\emph{ज्वालामुखी चाप} \omterm{def:b2:magmatic-series:series}{कैल्क-क्षारीय श्रेणियाँ} उगलते हैं। उनके \omterm{def:b2:elements-in-the-earth:trace}{अनुरेखी तत्त्व} एक
विशिष्ट चिह्न रखते हैं (नीचे का चित्र): \omterm{def:b2:elements-in-the-earth:bse}{आद्य प्रावार} पर प्रसामान्यीकृत करने पर चाप
बेसाल्ट Rb, Ba, Th, U, Pb और Sr में समृद्ध होते हैं, यानी उन तत्त्वों में जिन्हें स्लैब
प्रावार-कील में जोड़ता है, और अपने पड़ोसियों की तुलना में नायोबियम और टैंटलम में क्षीण होते हैं। आँकड़ा-समुच्चय का माध्यिका
चाप बेसाल्ट $\mathrm{Nb}/\mathrm{La} = 0.49$ रखता है, जबकि MORB के लिए यह 0.97
और महासागरीय-द्वीप बेसाल्टों के लिए 1.23 है।
% ledger: data:b4-basalt-settings

\begin{omfigure}
\centering
\begin{tikzpicture}
\begin{axis}[omrecord, width=0.86\linewidth, height=6.4cm, ymode=log,
  xmin=0.5, xmax=22.5, ymin=1, ymax=300,
  xtick={1,...,22}, xticklabels={Rb,Ba,Th,U,Nb,Ta,La,Ce,Pb,Sr,Nd,Zr,Hf,Sm,Eu,Ti,Gd,Dy,Y,Er,Yb,Lu},
  ylabel={नमूना / आद्य प्रावार},
  legend style={font=\scriptsize, at={(0.98,0.98)}, anchor=north east},
  tick label style={font=\scriptsize}, label style={font=\small}]
\addplot[thick, mark=square*, mark size=1.4pt, orange!85!black] table[x=idx, y=oib] {figdata/out/bachelor-2-series-spider.dat};
\addlegendentry{महासागरीय-द्वीप बेसाल्ट}
\addplot[thick, mark=*, mark size=1.4pt, blue!70!black] table[x=idx, y=morb] {figdata/out/bachelor-2-series-spider.dat};
\addlegendentry{MORB}
\addplot[thick, mark=triangle*, mark size=1.7pt, red!75!black] table[x=idx, y=iab] {figdata/out/bachelor-2-series-spider.dat};
\addlegendentry{द्वीप-चाप बेसाल्ट}
\end{axis}
\end{tikzpicture}
\par\smallskip{\small 241 MORB, 259
महासागरीय-द्वीप बेसाल्टों और 256 द्वीप-चाप बेसाल्टों के आद्य-प्रावार-प्रसामान्यीकृत माध्यिका संघटन (फ़र्मेश, 2006, का संकलन;
\omterm{def:b2:elements-in-the-earth:bse}{आद्य प्रावार}: मैकडोनो और सन, 1995)। चाप बेसाल्ट Nb
और Ta पर गर्त और Ba, Pb और Sr पर शिखर दिखाते हैं; महासागरीय-द्वीप बेसाल्ट सभी
सर्वाधिक \omterm{def:b2:elements-in-the-earth:partition}{असंगत तत्त्वों} में समृद्ध हैं।}
\end{omfigure}
% ledger: data:b4-basalt-settings, ms95:pm-trace

महाद्वीपों के भीतर बेसाल्टी \omterm{def:b2:magma-emplacement:magma}{मैग्मा} एक मोटी भूपर्पटी से मिलता है जिसे वह गरम कर सकता है और पिघला सकता है,
और परिणाम दो रूपों वाली ज्वालामुखीयता हो सकती है।

\begin{definition}[द्विरूपी ज्वालामुखीयता]\label{def:b2:magmatic-series:bimodal}
\emph{द्विरूपी ज्वालामुखीयता}\index{द्विरूपी ज्वालामुखीयता} किसी एक प्रांत में
मैग्माओं के दो समूहों, बेसाल्टी और रायोलाइटी, का उद्गार है, जिनके बीच
मध्यवर्ती संघटन बहुत कम होते हैं।
\end{definition}

येलोस्टोन हॉटस्पॉट के पथ के अनुदिश ज्वालामुखीयता द्विरूपी है: पूर्वी
स्नेक रिवर प्लेन पर लगभग \qty{75}{\percent} सिलिका वाले पुराने रायोलाइट
लगभग \qty{50}{\percent} वाले नए बेसाल्टों से ढके हैं। उनके बीच का अंतराल ही सुराग है:
अकेले किसी बेसाल्ट का विभेदन मध्यवर्ती लावा पीछे छोड़ता, और अंतराल
रायोलाइटों के किसी दूसरे स्रोत की ओर संकेत करता है, जैसे वह भूपर्पटी जिसे बेसाल्ट गरम करता है।
% ledger: yvo23:bimodal

\emph{बृहत आग्नेय प्रांत}, जिनसे पुस्तक 3 (स्नातक वर्ष 1) में भेंट हुई, इतने आयतनों में बेसाल्ट उगलते हैं
जिनकी बराबरी कोई अन्य परिवेश नहीं करता; उनकी श्रेणी उन्हीं आरेखों से पढ़ी जाती है।

\section{पृथक्करण आरेख}\label{sec:b2:magmatic-series:discrimination}

पिछले खंड के प्रतिरूप किसी विश्लेषण से सीधे परिवेश पढ़ने का सुझाव देते हैं,
और शैलिकीविदों ने ऐसा करने के लिए कई आरेख खींचे हैं।

\begin{definition}[पृथक्करण आरेख]\label{def:b2:magmatic-series:discrimination}
\emph{पृथक्करण आरेख}\index{पृथक्करण आरेख} रासायनिक चरों का, सामान्यतः
अचल \omterm{def:b2:elements-in-the-earth:trace}{अनुरेखी तत्त्वों} का, एक आरेख है, जो ज्ञात परिवेशों की चट्टानों के संघटनों के चारों ओर खींचे गए
क्षेत्रों में बँटा है, और जिसका उपयोग अज्ञात उद्गम वाली चट्टान को कोई परिवेश
सौंपने में होता है।
\end{definition}

क्लासिक आरेख पियर्स और कैन (1973) का Ti--Zr--Y त्रिभुज है, जो बेसाल्टों के लिए खींचा गया। यह
परिवर्तन से बहुत कम खिसकने वाले तत्त्वों का उपयोग करता है, जिससे इसे पुराने,
कायांतरित लावाओं पर लागू किया जा सकता है जिनका परिवेश खो चुका है। नीचे का चित्र इसे संकलन के उन 538 बेसाल्टों पर परखता है
जिनमें तीनों तत्त्व हैं। अधिकांश महासागरीय-द्वीप बेसाल्ट
अंतःप्लेट क्षेत्र D में पड़ते हैं, 173 में से 152, और अधिकांश MORB क्षेत्र B में, 164 में से 142; पर क्षेत्र B
चाप बेसाल्टों के साथ साझा है, और 201 चाप बेसाल्टों में से 101 भी वहीं पड़ते हैं।
% ledger: geoplotr:tizry, data:b4-basalt-settings

\begin{omfigure}
\centering
\begin{tikzpicture}
\begin{axis}[width=0.62\linewidth, height=0.55\linewidth, axis equal image,
  hide axis, xmin=-0.04, xmax=1.04, ymin=-0.06, ymax=0.92, clip=false,
  legend style={font=\scriptsize, at={(1.0,0.95)}, anchor=north west}]
\addplot[only marks, mark=*, mark size=0.7pt, blue!70!black, opacity=0.7] table[x=x, y=y] {figdata/out/bachelor-2-series-tizry-morb.dat};
\addlegendentry{MORB}
\addplot[only marks, mark=square*, mark size=0.7pt, orange!85!black, opacity=0.7] table[x=x, y=y] {figdata/out/bachelor-2-series-tizry-oib.dat};
\addlegendentry{महासागरीय-द्वीप बेसाल्ट}
\addplot[only marks, mark=triangle*, mark size=0.9pt, red!75!black, opacity=0.7] table[x=x, y=y] {figdata/out/bachelor-2-series-tizry-iab.dat};
\addlegendentry{द्वीप-चाप बेसाल्ट}
\foreach \f in {A,B,C,D}
  \addplot[thick, black, forget plot] table[x=x, y=y] {figdata/out/bachelor-2-series-tizry-\f.dat};
\draw[thick] (axis cs:0,0) -- (axis cs:1,0) -- (axis cs:0.5,0.8660254) -- cycle;
\node[anchor=south] at (axis cs:0.5,0.866) {Ti/100};
\node[anchor=north east] at (axis cs:0,0) {Zr};
\node[anchor=north west] at (axis cs:1,0) {Y $\times$ 3};
\draw[thin] (axis cs:0.33,0.36) -- (axis cs:0.22,0.46) node[anchor=south east, font=\small\bfseries, inner sep=1pt] {D};
\draw[thin] (axis cs:0.36,0.12) -- (axis cs:0.24,0.06) node[anchor=north east, font=\small\bfseries, inner sep=1pt] {C};
\draw[thin] (axis cs:0.56,0.37) -- (axis cs:0.76,0.50) node[anchor=west, font=\small\bfseries, inner sep=1pt] {A};
\draw[thin] (axis cs:0.53,0.21) -- (axis cs:0.72,0.13) node[anchor=north west, font=\small\bfseries, inner sep=1pt] {B};
\end{axis}
\end{tikzpicture}
\par\smallskip{\small ज्ञात परिवेश वाले बेसाल्टों के साथ पियर्स और कैन (1973) का Ti--Zr--Y
आरेख। A द्वीप-चाप थोलिआइट; B MORB, द्वीप-चाप थोलिआइट और
कैल्क-क्षारीय बेसाल्ट; C कैल्क-क्षारीय बेसाल्ट; D अंतःप्लेट बेसाल्ट। Ti
\ce{TiO2} से; तीनों ppm में।}
\end{omfigure}
% ledger: geoplotr:tizry, data:b4-basalt-settings

\begin{method}[पृथक्करण आरेख का उपयोग]\label{met:b2:magmatic-series:discrimination}
(1) जाँचिए कि चट्टान उसी प्रकार की है जिसके लिए आरेख बनाया गया था, सामान्यतः बेसाल्ट, और
कोई \omterm{def:b2:magmatic-differentiation:cumulate}{क्यूम्यूलेट} या विकसित लावा नहीं, जिसके \omterm{def:b2:elements-in-the-earth:trace}{अनुरेखी तत्त्व} स्रोत जितना ही
क्रिस्टलन को भी दर्शाते हैं। (2) परिवर्तित या कायांतरित चट्टानों के लिए परिवर्तन में अचल तत्त्वों (Ti, Zr, Y, Nb, भारी
दुर्लभ मृदाएँ) को वरीयता दीजिए। (3) नमूनों का एक समूह आलेखित कीजिए, एक नहीं:
किसी सीमा के निकट अकेला बिंदु कुछ तय नहीं करता। (4) परिणाम को एक संभाविता के रूप में पढ़िए:
क्षेत्र अपने अधिकांश संदर्भ नमूनों के चारों ओर खींचे गए थे, सभी के नहीं, और कुछ क्षेत्र
साझा हैं। (5) स्वतंत्र साक्ष्य से पुष्टि कीजिए: अन्य आरेख, समस्थानिक,
भूवैज्ञानिक संदर्भ।
\end{method}

चरण (4) को मात्रात्मक बनाया जा सकता है (\cref{exo:b2:magmatic-series:8})। यदि कोई बेसाल्ट क्षेत्र B में पड़ता है, और
तीनों परिवेश समान रूप से सामान्य हों, तो उसके MORB होने की संभावना केवल लगभग 0.62 है।
ऐसे संकलनों पर प्रशिक्षित सांख्यिकीय वर्गीकारक वही सीमा विरासत में पाते हैं: बेसाल्ट
अपने स्रोत और अपने गलन की स्मृति रखता है, और दो परिवेश समान
\omterm{def:b2:magma-emplacement:magma}{मैग्मा} बना सकते हैं।
% ledger: data:b4-basalt-settings

\begin{inthelab}[पूर्ण-चट्टान विश्लेषण की जाँच]
इनमें से किसी भी आरेख पर चट्टान को आलेखित करने से पहले शैलिकीविद विश्लेषण की जाँच करते हैं।
मुख्य ऑक्साइड, जो उदाहरण के लिए X-किरण प्रतिदीप्ति से मापे जाते हैं, वाष्पशीलों के साथ जोड़ने पर 100 के निकट
आने चाहिए: जल और कार्बन डाइऑक्साइड, जो सीधे मापे जाते हैं या साथ मिलकर
ज्वलन-हानि के रूप में, यानी वह द्रव्यमान जो चूर्ण तेज़ गरम करने पर खोता है। बड़ी ज्वलन-हानि
परिवर्तित चट्टान का चिह्न है, जिसके क्षार खिसक चुके हो सकते हैं; TAS
आरेख की पद्धति \qty{2}{\percent} से अधिक जल या
\qty{0.5}{\percent} से अधिक कार्बन डाइऑक्साइड वाले विश्लेषणों को अस्वीकार करती है। फिर विश्लेषणों को वाष्पशील-मुक्त 100 पर
पुनः परिकलित किया जाता है, लोहे को कुल \ce{FeO} के रूप में व्यक्त किया जाता है या फ़ेरस और फ़ेरिक में बाँटा जाता है,
और उसी दौर में विश्लेषित संदर्भ पदार्थों की तुलना उनके प्रमाणित
मानों से की जाती है (\cref{ch:b2:field-lab-methods-2})।
\end{inthelab}
% ledger: bgs:tas-use

\section{अभ्यास}

\begin{exercise}[$\star$]\label{exo:b2:magmatic-series:1}
TAS आरेख पर निम्नलिखित ज्वालामुखीय चट्टानों के नाम बताइए (\ce{SiO2}, \ce{Na2O}, \ce{K2O}
\unit{\percent} में, पहले से वाष्पशील-मुक्त): (क) 49.60, 2.22, 0.51 (हवाई का
संदर्भ बेसाल्ट BHVO-2); (ख) 58.5, 3.9, 1.6; (ग) 72.0, 4.0, 4.2; (घ) 47.0, 4.5, 2.5।
(घ) के लिए दोनों संभव चट्टानों के नाम बताइए और बताइए कि उनके बीच निर्णय किससे होता है।
\end{exercise}
% ledger: usgs:bhvo2-major, bgs:tas

\begin{exercise}[$\star$]\label{exo:b2:magmatic-series:2}
एक लावा का विश्लेषण: \ce{SiO2} 51.60, \ce{TiO2} 1.20, \ce{Al2O3} 16.10, \ce{FeO_T} 9.80,
\ce{MgO} 6.60, \ce{CaO} 8.20, \ce{Na2O} 3.50, \ce{K2O} 1.40, \ce{H2O} 1.40 और
\ce{CO2} \qty{0.20}{\percent}। (क) क्या यह विश्लेषण TAS आरेख के लिए स्वीकार्य है?
(ख) इसे वाष्पशील-मुक्त पुनः परिकलित कीजिए। (ग) पुनर्परिकलन से पहले और बाद में चट्टान का
नाम बताइए। (घ) क्या यह निम्न-, मध्यम- या उच्च-K है?
\end{exercise}
% ledger: bgs:tas-use, bgs:k2o-sio2

\begin{exercise}[$\star$]\label{exo:b2:magmatic-series:3}
\omterm{def:b2:magmatic-series:qapf}{QAPF आरेख} पर नाम बताइए: (क) \qty{25}{\percent} क्वार्ट्ज़,
\qty{30}{\percent} क्षार-फ़ेल्डस्पार, \qty{35}{\percent} प्लैजियोक्लेज़, \qty{8}{\percent}
बायोटाइट और \qty{2}{\percent} हॉर्नब्लेंड वाली एक स्थूल-कणिक चट्टान; (ख) \qty{5}{\percent} क्वार्ट्ज़,
\qty{5}{\percent} क्षार-फ़ेल्डस्पार, \qty{55}{\percent} प्लैजियोक्लेज़ और
\qty{35}{\percent} क्लाइनोपाइरॉक्सीन वाली एक चट्टान।
\end{exercise}
% ledger: bgs:qapf

\begin{exercise}[$\star$]\label{exo:b2:magmatic-series:4}
\ce{K2O}--\ce{SiO2} चित्र की रेखाओं से निम्न-, मध्यम- या उच्च-K में वर्गीकृत कीजिए:
(क) 50.0, 0.40; (ख) 56.0, 1.50; (ग) 62.0, 2.90; (घ) 66.0, 1.00 (\ce{SiO2},
\ce{K2O}, \unit{\percent})।
\end{exercise}
% ledger: bgs:k2o-sio2

\begin{exercise}[$\star\star$]\label{exo:b2:magmatic-series:5}
(क) $A$, $F$ और $M$ परिकलित कीजिए, \ce{FeO_T} 14.5, \ce{MgO} 4.5 और
\ce{Na2O + K2O} \qty{3.2}{\percent} वाले एक बेसाल्ट के लिए, और 6.1, 3.8 और 5.05 वाले एक एंडेसाइट के लिए।
(ख) विभाजक रेखा $(F, A, M) = (54.5, 23.0, 22.5)$, $(52.5, 30.0,
17.5)$ और $(49.6, 35.0, 15.4)$ से होकर जाती है, जिनमें पहला उसका सबसे लौह-समृद्ध बिंदु है;
थोलिआइटी ओर F की ओर है। हर चट्टान किस ओर पड़ती है? (ग) कौन-सा
खनिज, जल्दी क्रिस्टलित होकर, एंडेसाइट को वहीं रखता है जहाँ वह है?
\end{exercise}
% ledger: geoplotr:afm, vp21:series

\begin{exercise}[$\star\star$]\label{exo:b2:magmatic-series:6}
संकलन के माध्यिका MORB में \qty{4}{ppm} Nb और \qty{4.12}{ppm} La है,
माध्यिका महासागरीय-द्वीप बेसाल्ट में 38 और 30.8, माध्यिका चाप बेसाल्ट में 3 और 6.1। \omterm{def:b2:elements-in-the-earth:bse}{आद्य
प्रावार} के मानों Nb 0.658 और La \qty{0.648}{ppm} के साथ, हर एक के प्रसामान्यीकृत मान
और $\mathrm{Nb_N}/\mathrm{La_N}$ परिकलित कीजिए। किसमें नायोबियम विसंगति है, और क्यों?
\end{exercise}
% ledger: data:b4-basalt-settings, ms95:pm-trace

\begin{exercise}[$\star\star$]\label{exo:b2:magmatic-series:7}
Ti का द्रव्यमान-अंश $0.5993$ है, \ce{TiO2} में। तीन बेसाल्टों के Ti--Zr--Y आलेख-निर्देशांक
(Ti/100, Zr, $3\mathrm Y$, 100 पर पुनर्गठित) परिकलित कीजिए: (क) \ce{TiO2}
\qty{1.6}{\percent}, Zr 110, Y \qty{32}{ppm}; (ख) 2.8, 220, 29; (ग) 0.8, 50, 20। उन्हें
चित्र पर स्थित कीजिए और बताइए कि हर स्थिति से आप क्या निष्कर्ष निकाल सकते हैं।
\end{exercise}
% ledger: geoplotr:tizry

\begin{exercise}[$\star\star$]\label{exo:b2:magmatic-series:8}
इस अध्याय के Ti--Zr--Y परीक्षण में क्षेत्र B में 164 में से 142 MORB, 173 में से 4
महासागरीय-द्वीप बेसाल्ट और 201 में से 101 चाप बेसाल्ट हैं। (क) हर समूह का कितना अंश
B में पड़ता है? (ख) यदि तीनों परिवेश समान रूप से सामान्य हों, तो B में पड़ने वाले बेसाल्ट के MORB होने की
प्रायिकता क्या है? (ग) अकेला क्षेत्र B आपको क्या कहने देता है?
\end{exercise}
% ledger: data:b4-basalt-settings

\begin{exercise}[$\star\star$]\label{exo:b2:magmatic-series:9}
एक रायोलाइट में \ce{Na2O} 5.5, \ce{K2O} 4.5 और \ce{Al2O3} \qty{10.5}{\percent} है ($M$ =
61.98, 94.20 और \qty{101.96}{g/mol})। इसका आणविक $(\ce{Na2O +
K2O})/\ce{Al2O3}$ परिकलित कीजिए। (क) क्या यह अतिक्षारीय है? (ख) क्षार-फ़ेल्डस्पार प्रति \ce{Al2O3} एक
क्षार ऑक्साइड रखता है: क्षारों का कितना अंश उसमें प्रवेश नहीं कर सकता?
\end{exercise}
% ledger: bgs:tas-use

\begin{exercise}[$\star\star\star$]\label{exo:b2:magmatic-series:10}
लघुगणकीय-अनुपात प्रतिज्ञप्ति के मॉडल में $\lambda_A = 0.4$ और $\lambda_M =
0.8$ लीजिए। (क) $\lambda_F = 0.5$ और $\lambda_F =
0.7$ के लिए प्रवृत्ति की ढलान परिकलित कीजिए। (ख) आइसलैंडी चट्टानों की ढलान $-1.10$ है और कैस्केड की
$-0.72$: हर एक किस $\lambda_F$ का संकेत देती है? (ग) कौन-सी श्रेणी अपना लोहा अधिक तेज़ी से खोती है,
और कौन-सा खनिज इसे समझाता है? (घ) हर प्रवृत्ति के अनुदिश $F/M$ का क्या होता है?
\end{exercise}
% ledger: vp21:model

\begin{exercise}[$\star\star\star$]\label{exo:b2:magmatic-series:11}
एक सरलीकृत नॉर्म। ऑक्साइड के मोलों में काम कीजिए। सारा \ce{K2O} ऑर्थोक्लेज़
(\ce{K2O.Al2O3.6SiO2}) में डालिए, सारा \ce{Na2O} एल्बाइट (\ce{Na2O.Al2O3.6SiO2}) में,
शेष \ce{Al2O3} को उतने ही \ce{CaO} के साथ एनॉर्थाइट (\ce{CaO.Al2O3.2SiO2}) में,
शेष \ce{CaO} को उतने ही $(\ce{Mg,Fe})\ce{O}$ के साथ डायोप्साइड
(\ce{CaO.(Mg,Fe)O.2SiO2}) में, और बचे हुए $(\ce{Mg,Fe})\ce{O}$ को हाइपरस्थीन
(\ce{(Mg,Fe)O.SiO2}) में। यदि सिलिका कम पड़े, तो हाइपरस्थीन को ओलिवीन
(\ce{2(Mg,Fe)O.SiO2}) में बदलिए; यदि फिर भी कम पड़े, तो एल्बाइट को नेफ़लीन
(\ce{Na2O.Al2O3.2SiO2}) में बदलिए। इसे लागू कीजिए (क) BHVO-2 पर: \ce{SiO2} 49.60, \ce{Al2O3} 13.44,
\ce{FeO_T} 11.15, \ce{MgO} 7.257, \ce{CaO} 11.40, \ce{Na2O} 2.219, \ce{K2O}
\qty{0.513}{\percent}; (ख) 45.5, 15.0, 11.5, 8.5, 10.0, 3.8 और 1.5 वाले एक बेसाल्ट पर (मोलर
द्रव्यमान \ce{SiO2} 60.08, \ce{Al2O3} 101.96, \ce{FeO} 71.84, \ce{MgO} 40.30, \ce{CaO}
56.08, \ce{Na2O} 61.98, \ce{K2O} \qty{94.20}{g/mol})। क्या हर एक सिलिका-संतृप्त है?
हर बेसाल्ट किस श्रेणी का है?
\end{exercise}
% ledger: usgs:bhvo2-major, bgs:series-terms

\begin{exercise}[$\star\star\star$]\label{exo:b2:magmatic-series:12}
संकलन में $\mathrm{Nb}/\mathrm{La} < 0.7$, दोनों तत्त्वों के लिए विश्लेषित 141 में से 116 चाप बेसाल्टों,
101 में से 21 MORB और 137 में से 3 महासागरीय-द्वीप बेसाल्टों के लिए सत्य है। (क) तीनों
अंश परिकलित कीजिए। (ख) तीनों परिवेशों के समान अनुपातों के साथ, $\mathrm{Nb}/\mathrm{La} < 0.7$ वाले बेसाल्ट के
चाप बेसाल्ट होने की प्रायिकता क्या है? (ग)
इसे ऐसे क्षेत्र के लिए दोबारा कीजिए जहाँ \qty{60}{\percent} बेसाल्ट MORB हैं, \qty{30}{\percent}
महासागरीय-द्वीप बेसाल्ट और \qty{10}{\percent} चाप बेसाल्ट। (घ) भाग (ग) पृथक्करण आरेखों के उपयोग के लिए
क्या सबक देता है?
\end{exercise}
% ledger: data:b4-basalt-settings

\section{समस्या: एक ज्वालामुखी, तीन लावा}

\begin{problem}[{सप्ताहांत समस्या --- तीन विश्लेषणों से एक ज्वालामुखी के परिवेश तक और
उन क्रिस्टलों तक जो वह पीछे छोड़ गया}]
\label{pb:b2:magmatic-series:1}
एक (काल्पनिक) स्तरित ज्वालामुखी के तीन लावा, उसके सबसे पुराने पार्श्व प्रवाहों से उसके सबसे नए
गुंबद तक, ये मान देते हैं (\unit{\percent}):

\smallskip\noindent{\small
\begin{tabular}{@{}lrrrrrrrrrr@{}}
 & \ce{SiO2} & \ce{TiO2} & \ce{Al2O3} & \ce{FeO_T} & \ce{MgO} & \ce{CaO} & \ce{Na2O} & \ce{K2O} & \ce{H2O} & \ce{CO2}\\
लावा 1 & 50.4 & 0.95 & 17.6 & 8.6 & 7.9 & 10.2 & 2.9 & 0.62 & 0.70 & 0.10\\
लावा 2 & 58.6 & 0.80 & 17.4 & 6.1 & 3.8 & 6.9 & 3.9 & 1.15 & 1.10 & 0.05\\
लावा 3 & 65.3 & 0.55 & 16.2 & 4.0 & 1.6 & 4.3 & 4.4 & 1.85 & 1.20 & 0.05\\
\end{tabular}}

\smallskip\noindent और ppm में: Rb 10, 20, 32; Th 1.1, 2.1, 3.4; Nb 2.0, 3.6, 5.6;
La 7, 13, 20; Sr 520, 480, 380; Zr 70, 125, 190; Y 18, 22, 24। \omterm{def:b2:elements-in-the-earth:bse}{आद्य प्रावार}: Nb
0.658, La 0.648, Th \qty{0.0795}{ppm}।

\noindent\textbf{भाग I --- लावाओं का नामकरण।}
\begin{enumerate}
\item हर विश्लेषण का योग कीजिए। क्या तीनों TAS आरेख के लिए स्वीकार्य हैं?
\item लावा 1 को वाष्पशील-मुक्त पुनः परिकलित कीजिए: गुणक, \ce{SiO2} और
\ce{Na2O + K2O} बताइए।
\item लावा 2 और 3 के लिए भी यही कीजिए।
\item TAS आरेख पर तीनों लावाओं के नाम बताइए।
\item क्या पुनर्परिकलन के बिना नाम बदल जाते?
\item लावा 2 के संघटन वाली एक स्थूल-कणिक चट्टान ज्वालामुखी की
जड़ों में क्रिस्टलित हुई। कौन-सी पद्धति उसका नाम रखती है, और आप उसे किस क्षेत्र में अपेक्षित करेंगे?
\end{enumerate}

\noindent\textbf{भाग II --- कौन-सी श्रेणी?}
\begin{enumerate}[resume]
\item हर लावा को निम्न-, मध्यम- या उच्च-K में वर्गीकृत कीजिए।
\item हर एक के लिए $A$, $F$, $M$ परिकलित कीजिए (वाष्पशील-मुक्त)।
\item विभाजक रेखा $(F, A, M) = (42.8, 17.1, 40.1)$, $(46.9, 18.1,
35.0)$, $(52.5, 30.0, 17.5)$, $(49.6, 35.0, 15.4)$, $(40.8, 50.0, 9.2)$ और $(37.9,
55.0, 7.1)$ से होकर जाती है, और थोलिआइटी क्षेत्र F की ओर है। पड़ोसी बिंदुओं के बीच $F$ को
$A$ में रैखिक रूप से अंतर्वेशित करते हुए, हर लावा किस ओर पड़ता है?
\item हर लावा के लिए \ce{FeO_T}/\ce{MgO} परिकलित कीजिए। क्या लोहे का संवर्धन है?
\item ज्वालामुखी किस श्रेणी का है? कौन-सा खनिज और कौन-सी परिस्थितियाँ
इसे समझाती हैं?
\end{enumerate}

\noindent\textbf{भाग III --- कौन-सा परिवेश?}
\begin{enumerate}[resume]
\item लावा 1 के Nb, La और Th को \omterm{def:b2:elements-in-the-earth:bse}{आद्य प्रावार} पर प्रसामान्यीकृत कीजिए।
\item लावा 1 के लिए $\mathrm{Nb_N}/\mathrm{La_N}$ परिकलित कीजिए और उसकी तुलना
\cref{exo:b2:magmatic-series:6} की माध्यिकाओं से कीजिए।
\item $\mathrm{Th_N}/\mathrm{Nb_N}$ परिकलित कीजिए। थोरियम को कौन लाया?
\item लावा 1 के Ti--Zr--Y निर्देशांक परिकलित कीजिए (वाष्पशील-मुक्त
विश्लेषण का \ce{TiO2})। उसका क्षेत्र क्या कहता है? लावा 2 और 3 को आलेखित क्यों न करें?
\item निष्कर्ष निकालिए: यह ज्वालामुखी कहाँ बढ़ा?
\end{enumerate}

\noindent\textbf{भाग IV --- कितना क्रिस्टलित हुआ?}
\begin{enumerate}[resume]
\item पोटाश आरंभिक खनिजों में से किसी में प्रवेश नहीं करता। लावा 1 और 2 के वाष्पशील-मुक्त \ce{K2O} से,
लावा 2 के बनने के समय बचे द्रव का अंश $F$ परिकलित कीजिए, यदि वह
लावा 1 से \omterm{def:b2:magmatic-differentiation:crystallisation}{भिन्नात्मक क्रिस्टलन} द्वारा व्युत्पन्न है।
\item रुबिडियम और थोरियम के साथ भी यही कीजिए।
\item ज़िरकोनियम का समग्र $D \approx 0.1$ है। उससे $F$ का अनुमान लगाने के लिए रैले नियम का
उपयोग कीजिए।
\item लावा 2 से लावा 3 तक स्ट्रॉन्शियम घटता है। लावा 2 और 3 के बीच पोटाश से $F$ परिकलित कीजिए,
और उसका उपयोग उस अंतराल पर Sr के समग्र $D$ का अनुमान लगाने में कीजिए। कौन-सा खनिज
उत्तरदायी है?
\item लावा 1 से लावा 3 तक द्रव का कितना अंश क्रिस्टलित हुआ?
\item यदि \omterm{def:b2:magma-emplacement:magma}{मैग्मा} ने पोटाश से समृद्ध ग्रेनाइटी भूपर्पटी आत्मसात की होती, तो क्या पोटाश वाला
अनुमान बहुत ऊँचा होता या बहुत नीचा?
\item लावा 1 से 2 के लिए भाग IV के चारों अनुमानों का औसत लीजिए और एंडेसाइट बनाने के लिए
बेसाल्ट के हर \qty{100}{kg} पर पीछे छूटे क्रिस्टलों का द्रव्यमान बताइए, जो इस समस्या की
मुख्य संख्या है।
\end{enumerate}
\end{problem}
% ledger: bgs:tas-use, bgs:tas, bgs:k2o-sio2, geoplotr:afm, geoplotr:tizry, ms95:pm-trace
